\begin{minipage}{0.48\linewidth}
    Un skieur, de masse $M=\qty{80}{\kg}$ (avec son équipement), se trouve au
    sommet d'une piste rectiligne de pente $8 \%$ (la piste descend
    verticalement de $\qty{8}{\m}$ pour un parcours de $\qty{100}{\m}$ le long
    de la piste). La longueur de la pente est $L=\qty{200}{\m}$.
\end{minipage}
\hfill
\begin{minipage}{0.44\linewidth}
    \begin{figure}[H]
        \centering
        \includegraphics[width=\textwidth]{2026-03-20_183711-}
    \end{figure}
\end{minipage}
\begin{enumerate}
    \item Le skieur part du sommet de la piste avec une vitesse de valeur
          $v=\qty{0.5}{\m\per\s}$.
    
          Calculer, au moment du départ, son énergie cinétique et son énergie
          potentielle de pesanteur.
    \item Lorsque le skieur descend, quelle transformation d'énergie se
          produit-il~?
    \item Si la neige est verglacée (absence de frottements) et si la résistance
          de l'air est négligeable, comment sont liées les deux formes d'énergie
          que possède le skieur~?
    \item Combien vaut l'énergie cinétique du skieur lorsqu'il est parvenu au
          bas de la piste (point $P$).

          En déduire la valeur de sa vitesse au point $P$.
    \item Que peut-on dire de la valeur de sa vitesse s'il poursuit son chemin
          sur la piste horizontale~? Justifier la réponse.
\end{enumerate}

